% =============================================================================
%  progress-report-04.tex — รายงานความก้าวหน้าโครงการ ULMs ครั้งที่ 4
%  *** ฉบับส่งอาจารย์ รอบ 11 – 17 ก.ย. 2569 (วันที่ 7 ของ Sprint 4) ***
%
%  build:  cd ~/Projects/SoftwareEn/docs && latexmk progress-report-04.tex
%          ผลลัพธ์ออกที่ build/progress-report-04.pdf
%
%  หลักการเขียน — ยกมาจาก progress-report-03.tex ทุกข้อ อ่านหัวไฟล์นั้นก่อนแก้
%    สั้น ไม่มีสารบัญ · แผนเส้นเดียว · อ้างได้เฉพาะฉบับที่ส่งจริง (ครั้งที่ 1 ถึง 3)
%    ห้ามรายงานเป็นจำนวนบรรทัด · ไม่มี \textbf กลางย่อหน้า
%    ศัพท์ของการพัฒนาใช้อังกฤษ และเว้นวรรคหน้าหลังทุกครั้งที่ชนอักษรไทย
%    ห้ามรัน regex ยุบช่องไฟทับทั้งไฟล์
%
%  ห้ามแก้ progress-report-01/02/03.tex ย้อนหลัง ทั้งสามฉบับส่งไปแล้ว
%
%  ตัวเลขทุกตัวอ้าง origin/main ที่ 0cb21a7 (17 ก.ย. 2569) fetch ซ้ำก่อนสร้าง PDF แล้ว
%  ไม่นับ spike-option-b/ และไม่นับ branch ที่ยังไม่ merge
%
%  วันรายงานรอบนี้ไม่ตรงขอบ Sprint — เป็นวันที่ 7 จาก 14 ของ Sprint 4 พอดี
%  ค่าแผนจึงสอดแทรกครึ่งทางระหว่างแถว S3 กับ S4 ของภาคผนวก เหมือนที่ครั้งที่ 2 ทำ
%
%  ที่มาของตัวเลข ตรวจซ้ำได้จาก repo:
%    86 procedure / 10 domain = ผลรวม @Query/@Mutation ใน backend/src/*/*.router.ts
%                    เท่ากับรอบก่อนพอดี (นับที่ c8ed505 ได้ 86 เท่ากัน)
%    30 ตาราง 9 enum 10 migration = backend/prisma/ (migration เพิ่ม 1 คือ timestamptz)
%    266 test / 24 suite = `npm test` ในเครื่องนี้ ชี้ DATABASE_URL ไปฐาน ulms_test
%                    ผ่านทั้งหมด (รอบก่อนรายงาน 234 เพราะยังไม่ได้ต่อ Postgres)
%    10 จาก 23 service มี .spec.ts คู่กัน (รอบก่อน 8 จาก 21 วิธีนับเดียวกัน)
%    33 vitest / 3 ไฟล์ · 49 AC check · lint 0 error 7 warning
%    12 ผ่าน 1 ข้าม จาก 13 = `npx playwright test` โดยสตาร์ต backend :3000 และ
%                    frontend :5173 เองบนฐาน ulms_e2e ที่ seed ใหม่
%    29 หน้า = frontend/src/features/**/*-page.tsx (รอบก่อน 28 เพิ่มหน้า easter egg)
%    27 มี UI จริง = ไม่ใช้ PlaceholderPage (เหลือ handover กับ reports/export)
%    5 หน้าอยู่ในโดเมนที่ไม่มี backend เลย = ห้อง 3 หน้า อุทธรณ์ 2 หน้า (เท่ารอบก่อน)
%    12 หน้ายัง import mock-data.ts แต่ 8 ใน 12 เป็นค่าคงที่ร่วม ไม่ใช่ข้อมูลจำลอง
%    77 procedure / 11 domain ที่ FE ประกาศใน frontend/src/server/trpc-contract.ts
%    62 จาก 65 รหัสข้อผิดพลาดมีคำแปลไทย (รอบก่อน 60 จาก 63)
%    29 commit / 10 PR ในช่วง c8ed505..0cb21a7 (ไม่รวม merge = 19)
% =============================================================================
% =============================================================================
%  ulms-preamble.tex — preamble ที่ใช้ร่วมกันทุกเอกสารของโครงการ ULMs
%
%  ใช้โดย:  main.tex (ข้อเสนอโครงการ) · progress.tex (รายงานความก้าวหน้า)
%
%  เอกสารที่เรียกใช้ต้องทำสองอย่างหลัง % =============================================================================
%  ulms-preamble.tex — preamble ที่ใช้ร่วมกันทุกเอกสารของโครงการ ULMs
%
%  ใช้โดย:  main.tex (ข้อเสนอโครงการ) · progress.tex (รายงานความก้าวหน้า)
%
%  เอกสารที่เรียกใช้ต้องทำสองอย่างหลัง % =============================================================================
%  ulms-preamble.tex — preamble ที่ใช้ร่วมกันทุกเอกสารของโครงการ ULMs
%
%  ใช้โดย:  main.tex (ข้อเสนอโครงการ) · progress.tex (รายงานความก้าวหน้า)
%
%  เอกสารที่เรียกใช้ต้องทำสองอย่างหลัง \input{ulms-preamble}:
%    1) \hypersetup{pdftitle={...}}  — ชื่อเอกสารใน metadata ของ PDF
%    2) \begin{document}
%  (pdftitle ถูกย้ายออกจากไฟล์นี้ เพราะเป็นค่าเฉพาะของแต่ละเอกสาร)
%
%  แก้ที่นี่ที่เดียว = มีผลกับทุกเอกสาร ห้ามก๊อป preamble ไปวางซ้ำ
% =============================================================================
% =============================================================================
%  main.tex  —  เอกสารข้อเสนอโครงการ (Academic Project Proposal)
%  ระบบบริหารจัดการการยืม–คืนอุปกรณ์มหาวิทยาลัย (ULMs)
%
%  IMPORTANT: This document contains Thai text and MUST be compiled with
%  XeLaTeX (not pdflatex). Thai needs a Unicode font + ICU line-breaking.
%      Build:  latexmk main.tex        (uses .latexmkrc -> xelatex)
%      or:     xelatex main.tex        (run twice for the table of contents)
% =============================================================================

\documentclass[11pt, a4paper]{article}

% ---- Fonts & Thai support ----------------------------------------------------
% fontspec lets XeLaTeX use any system font. Laksaman is a complete serif face
% that covers BOTH Thai and Latin glyphs, so one \setmainfont handles the whole
% (Thai + English) document without font switching.
\usepackage{fontspec}

% Thai has no spaces between words, so TeX cannot find line-break points on its
% own. These two XeTeX primitives switch on ICU's dictionary-based Thai word
% segmentation, which is what makes Thai paragraphs wrap correctly.
\XeTeXlinebreaklocale "th"
\XeTeXlinebreakskip = 0pt plus 0.1pt

\setmainfont{Laksaman}

% ---- Page geometry & paragraph style ----------------------------------------
\usepackage[margin=2.5cm]{geometry}
\usepackage{parskip}          % block paragraphs (no indent, gap between) — common
                              % and more readable for Thai documents
\linespread{1.15}             % a little extra leading; Thai has tall tone marks

% ---- Colours (monochrome / grayscale palette) ------------------------------
\usepackage[table]{xcolor}    % 'table' option enables \rowcolor for table headers
% Monochrome (grayscale) palette. Names kept as-is so every reference still
% works; only the tones changed from the original green/blue accents.
\definecolor{kugreen}{HTML}{1A1A1A}   % near-black — section headings & title
\definecolor{kudark}{HTML}{333333}    % dark gray — subsection/subsubsection titles
\definecolor{headbg}{HTML}{EAEAEA}    % light gray background for table header rows
\definecolor{linkblue}{HTML}{404040}  % medium-dark gray — links / citations

% ---- Section heading styling -------------------------------------------------
% 'nobottomtitles*' stops a heading from being printed near the bottom of a
% page: if less than \bottomtitlespace of room is left, the heading (and its
% following text) moves to the next page. This is titlesec's built-in, robust
% replacement for the earlier \needspace hack — the latter injected vertical
% glue that clashed with longtable page-splitting ("Infinite glue shrinkage").
\usepackage[nobottomtitles*]{titlesec}
\titleformat{\section}
  {\color{kugreen}\Large\bfseries}{\thesection}{0.6em}{}
\titleformat{\subsection}
  {\color{kudark}\large\bfseries}{\thesubsection}{0.5em}{}
\titleformat{\subsubsection}
  {\color{kudark}\normalsize\bfseries}{\thesubsubsection}{0.5em}{}
\titlespacing*{\section}{0pt}{1.4\baselineskip}{0.5\baselineskip}

% Minimum room a heading needs at the bottom of a page; if only ~1–3 lines are
% left (less than this), the heading is bumped to the next page instead of being
% stranded above its table/text.
\renewcommand{\bottomtitlespace}{4\baselineskip}

% ---- Hierarchical indentation by heading level ------------------------------
% Each heading level shifts BOTH the heading and the body text under it further
% to the right: section = 0, subsection (x.y) = 1 step, subsubsection (x.y.z) =
% 2 steps. We do this by having the sectioning commands set \leftskip (a global
% left margin that carries through every following line until the next heading
% resets it). \par first, so the *previous* paragraph is not re-indented.
% \leftskip indents plain paragraphs; \@totalleftmargin makes list environments
% (itemize/enumerate) start from the same indent instead of the page margin.
\makeatletter
\newlength{\lvlindent}
\setlength{\lvlindent}{1.8em}          % one indent step
\let\ulmsSecIndent\section
\renewcommand{\section}{\par\global\leftskip=0pt\global\@totalleftmargin=0pt\relax\ulmsSecIndent}
\let\ulmsSubsecIndent\subsection
\renewcommand{\subsection}{\par\global\leftskip=\lvlindent\global\@totalleftmargin=\lvlindent\relax\ulmsSubsecIndent}
\let\ulmsSubsubsecIndent\subsubsection
\renewcommand{\subsubsection}{\par\global\leftskip=2\lvlindent\global\@totalleftmargin=2\lvlindent\relax\ulmsSubsubsecIndent}
\makeatother

% Indent the headings themselves to match (titlesec ignores \leftskip for the
% heading line, so its left margin is set here via the first \titlespacing arg).
\titlespacing*{\subsection}   {\lvlindent} {1.0\baselineskip}{0.4\baselineskip}
\titlespacing*{\subsubsection}{2\lvlindent}{0.8\baselineskip}{0.3\baselineskip}

% ---- Lists (tighter, custom bullets) ----------------------------------------
\usepackage{enumitem}
\setlist[itemize]{leftmargin=1.6em, itemsep=2pt, topsep=3pt}
\setlist[enumerate]{leftmargin=1.9em, itemsep=2pt, topsep=3pt}

% ---- Tables ------------------------------------------------------------------
\usepackage{array}
\usepackage{booktabs}         % \toprule / \midrule / \bottomrule
\usepackage{tabularx}         % auto-width columns that wrap long Thai text
\usepackage{longtable}        % tables that can break across pages
\usepackage{multirow}         % merged cells for the team-structure table

% Left-aligned wrapping column types:
%   L{width} — fixed-width ragged-right paragraph column
%   Y        — flexible ragged-right X column (for tabularx)
\newcolumntype{L}[1]{>{\raggedright\arraybackslash}p{#1}}
\newcolumntype{Y}{>{\raggedright\arraybackslash}X}
\renewcommand{\arraystretch}{1.3}

% Helpers for a styled header row inside tables.
%   \hrow  — starts a header row and shades it (must be FIRST in the row;
%            colortbl allows only one \rowcolor per row).
%   \thead — bold header-cell text.
\newcommand{\hrow}{\rowcolor{headbg}}
\newcommand{\thead}[1]{\textbf{#1}}

% \begin{keeptable} … \end{keeptable} — keeps a table and the bold label above
% it on the SAME page (used by the data dictionary in §5). A tabularx cannot
% split across pages, so LaTeX is free to break between the label and its
% table, stranding the label alone at the bottom of a page; wrapping the pair
% in a minipage makes them one unbreakable block.
%   * \leftskip is zeroed inside: the minipage box is already placed at the
%     current heading indent, so leaving it set would indent the content twice
%     and push the table past the right margin.
%   * \parskip is restored because minipage's \@parboxrestore zeroes it, which
%     would glue the table straight onto the label line.
% Tables inside must use \linewidth (= the minipage width), not \textwidth.
\newenvironment{keeptable}
  {\par\medskip\noindent
   \begin{minipage}{\dimexpr\textwidth-\leftskip\relax}%
   \leftskip=0pt\relax
   \setlength{\parskip}{0.5\baselineskip plus 2pt}}
  {\end{minipage}\par\medskip}

% ---- Table of contents styling ----------------------------------------------
% By default LaTeX gives dotted leaders to subsections but NOT to sections
% (they print bold with just a right-aligned page number). tocloft lets us turn
% on dots for the section level too, so every TOC line runs "….." to its page.
\usepackage{tocloft}
\renewcommand{\cftsecdotsep}{\cftdotsep}   % dotted leaders for section entries
\renewcommand{\cftsecleader}{\cftdotfill{\cftsecdotsep}}
\renewcommand{\cftsecpagefont}{\normalfont}  % keep page numbers un-bold

% ---- Figures -----------------------------------------------------------------
\usepackage{graphicx}
\graphicspath{{figures/}}

% Landscape (rotated) full-page floats — needed by the ER diagram in §5, which
% is far wider than it is tall and would be unreadable scaled to \textwidth.
\usepackage{rotating}

% ---- Flowcharts / diagrams (TikZ) -------------------------------------------
% Used for the borrow–return workflow diagram (§5). Libraries: geometric shapes
% for decision diamonds, arrows.meta for arrowheads, positioning for relative
% node placement, calc/fit/backgrounds for routing loops and grouping.
\usepackage{tikz}
\usetikzlibrary{shapes.geometric, arrows.meta, positioning, calc, fit, backgrounds}

% ---- Links & bookmarks -------------------------------------------------------
\usepackage{hyperref}
\hypersetup{
  colorlinks=true, linkcolor=black, citecolor=linkblue, urlcolor=linkblue,
  pdflang={th},
}

% Thai labels for auto-generated headings.
\renewcommand{\contentsname}{สารบัญ (Table of Contents)}
\renewcommand{\tablename}{ตาราง}
\renewcommand{\figurename}{รูปที่}

:
%    1) \hypersetup{pdftitle={...}}  — ชื่อเอกสารใน metadata ของ PDF
%    2) \begin{document}
%  (pdftitle ถูกย้ายออกจากไฟล์นี้ เพราะเป็นค่าเฉพาะของแต่ละเอกสาร)
%
%  แก้ที่นี่ที่เดียว = มีผลกับทุกเอกสาร ห้ามก๊อป preamble ไปวางซ้ำ
% =============================================================================
% =============================================================================
%  main.tex  —  เอกสารข้อเสนอโครงการ (Academic Project Proposal)
%  ระบบบริหารจัดการการยืม–คืนอุปกรณ์มหาวิทยาลัย (ULMs)
%
%  IMPORTANT: This document contains Thai text and MUST be compiled with
%  XeLaTeX (not pdflatex). Thai needs a Unicode font + ICU line-breaking.
%      Build:  latexmk main.tex        (uses .latexmkrc -> xelatex)
%      or:     xelatex main.tex        (run twice for the table of contents)
% =============================================================================

\documentclass[11pt, a4paper]{article}

% ---- Fonts & Thai support ----------------------------------------------------
% fontspec lets XeLaTeX use any system font. Laksaman is a complete serif face
% that covers BOTH Thai and Latin glyphs, so one \setmainfont handles the whole
% (Thai + English) document without font switching.
\usepackage{fontspec}

% Thai has no spaces between words, so TeX cannot find line-break points on its
% own. These two XeTeX primitives switch on ICU's dictionary-based Thai word
% segmentation, which is what makes Thai paragraphs wrap correctly.
\XeTeXlinebreaklocale "th"
\XeTeXlinebreakskip = 0pt plus 0.1pt

\setmainfont{Laksaman}

% ---- Page geometry & paragraph style ----------------------------------------
\usepackage[margin=2.5cm]{geometry}
\usepackage{parskip}          % block paragraphs (no indent, gap between) — common
                              % and more readable for Thai documents
\linespread{1.15}             % a little extra leading; Thai has tall tone marks

% ---- Colours (monochrome / grayscale palette) ------------------------------
\usepackage[table]{xcolor}    % 'table' option enables \rowcolor for table headers
% Monochrome (grayscale) palette. Names kept as-is so every reference still
% works; only the tones changed from the original green/blue accents.
\definecolor{kugreen}{HTML}{1A1A1A}   % near-black — section headings & title
\definecolor{kudark}{HTML}{333333}    % dark gray — subsection/subsubsection titles
\definecolor{headbg}{HTML}{EAEAEA}    % light gray background for table header rows
\definecolor{linkblue}{HTML}{404040}  % medium-dark gray — links / citations

% ---- Section heading styling -------------------------------------------------
% 'nobottomtitles*' stops a heading from being printed near the bottom of a
% page: if less than \bottomtitlespace of room is left, the heading (and its
% following text) moves to the next page. This is titlesec's built-in, robust
% replacement for the earlier \needspace hack — the latter injected vertical
% glue that clashed with longtable page-splitting ("Infinite glue shrinkage").
\usepackage[nobottomtitles*]{titlesec}
\titleformat{\section}
  {\color{kugreen}\Large\bfseries}{\thesection}{0.6em}{}
\titleformat{\subsection}
  {\color{kudark}\large\bfseries}{\thesubsection}{0.5em}{}
\titleformat{\subsubsection}
  {\color{kudark}\normalsize\bfseries}{\thesubsubsection}{0.5em}{}
\titlespacing*{\section}{0pt}{1.4\baselineskip}{0.5\baselineskip}

% Minimum room a heading needs at the bottom of a page; if only ~1–3 lines are
% left (less than this), the heading is bumped to the next page instead of being
% stranded above its table/text.
\renewcommand{\bottomtitlespace}{4\baselineskip}

% ---- Hierarchical indentation by heading level ------------------------------
% Each heading level shifts BOTH the heading and the body text under it further
% to the right: section = 0, subsection (x.y) = 1 step, subsubsection (x.y.z) =
% 2 steps. We do this by having the sectioning commands set \leftskip (a global
% left margin that carries through every following line until the next heading
% resets it). \par first, so the *previous* paragraph is not re-indented.
% \leftskip indents plain paragraphs; \@totalleftmargin makes list environments
% (itemize/enumerate) start from the same indent instead of the page margin.
\makeatletter
\newlength{\lvlindent}
\setlength{\lvlindent}{1.8em}          % one indent step
\let\ulmsSecIndent\section
\renewcommand{\section}{\par\global\leftskip=0pt\global\@totalleftmargin=0pt\relax\ulmsSecIndent}
\let\ulmsSubsecIndent\subsection
\renewcommand{\subsection}{\par\global\leftskip=\lvlindent\global\@totalleftmargin=\lvlindent\relax\ulmsSubsecIndent}
\let\ulmsSubsubsecIndent\subsubsection
\renewcommand{\subsubsection}{\par\global\leftskip=2\lvlindent\global\@totalleftmargin=2\lvlindent\relax\ulmsSubsubsecIndent}
\makeatother

% Indent the headings themselves to match (titlesec ignores \leftskip for the
% heading line, so its left margin is set here via the first \titlespacing arg).
\titlespacing*{\subsection}   {\lvlindent} {1.0\baselineskip}{0.4\baselineskip}
\titlespacing*{\subsubsection}{2\lvlindent}{0.8\baselineskip}{0.3\baselineskip}

% ---- Lists (tighter, custom bullets) ----------------------------------------
\usepackage{enumitem}
\setlist[itemize]{leftmargin=1.6em, itemsep=2pt, topsep=3pt}
\setlist[enumerate]{leftmargin=1.9em, itemsep=2pt, topsep=3pt}

% ---- Tables ------------------------------------------------------------------
\usepackage{array}
\usepackage{booktabs}         % \toprule / \midrule / \bottomrule
\usepackage{tabularx}         % auto-width columns that wrap long Thai text
\usepackage{longtable}        % tables that can break across pages
\usepackage{multirow}         % merged cells for the team-structure table

% Left-aligned wrapping column types:
%   L{width} — fixed-width ragged-right paragraph column
%   Y        — flexible ragged-right X column (for tabularx)
\newcolumntype{L}[1]{>{\raggedright\arraybackslash}p{#1}}
\newcolumntype{Y}{>{\raggedright\arraybackslash}X}
\renewcommand{\arraystretch}{1.3}

% Helpers for a styled header row inside tables.
%   \hrow  — starts a header row and shades it (must be FIRST in the row;
%            colortbl allows only one \rowcolor per row).
%   \thead — bold header-cell text.
\newcommand{\hrow}{\rowcolor{headbg}}
\newcommand{\thead}[1]{\textbf{#1}}

% \begin{keeptable} … \end{keeptable} — keeps a table and the bold label above
% it on the SAME page (used by the data dictionary in §5). A tabularx cannot
% split across pages, so LaTeX is free to break between the label and its
% table, stranding the label alone at the bottom of a page; wrapping the pair
% in a minipage makes them one unbreakable block.
%   * \leftskip is zeroed inside: the minipage box is already placed at the
%     current heading indent, so leaving it set would indent the content twice
%     and push the table past the right margin.
%   * \parskip is restored because minipage's \@parboxrestore zeroes it, which
%     would glue the table straight onto the label line.
% Tables inside must use \linewidth (= the minipage width), not \textwidth.
\newenvironment{keeptable}
  {\par\medskip\noindent
   \begin{minipage}{\dimexpr\textwidth-\leftskip\relax}%
   \leftskip=0pt\relax
   \setlength{\parskip}{0.5\baselineskip plus 2pt}}
  {\end{minipage}\par\medskip}

% ---- Table of contents styling ----------------------------------------------
% By default LaTeX gives dotted leaders to subsections but NOT to sections
% (they print bold with just a right-aligned page number). tocloft lets us turn
% on dots for the section level too, so every TOC line runs "….." to its page.
\usepackage{tocloft}
\renewcommand{\cftsecdotsep}{\cftdotsep}   % dotted leaders for section entries
\renewcommand{\cftsecleader}{\cftdotfill{\cftsecdotsep}}
\renewcommand{\cftsecpagefont}{\normalfont}  % keep page numbers un-bold

% ---- Figures -----------------------------------------------------------------
\usepackage{graphicx}
\graphicspath{{figures/}}

% Landscape (rotated) full-page floats — needed by the ER diagram in §5, which
% is far wider than it is tall and would be unreadable scaled to \textwidth.
\usepackage{rotating}

% ---- Flowcharts / diagrams (TikZ) -------------------------------------------
% Used for the borrow–return workflow diagram (§5). Libraries: geometric shapes
% for decision diamonds, arrows.meta for arrowheads, positioning for relative
% node placement, calc/fit/backgrounds for routing loops and grouping.
\usepackage{tikz}
\usetikzlibrary{shapes.geometric, arrows.meta, positioning, calc, fit, backgrounds}

% ---- Links & bookmarks -------------------------------------------------------
\usepackage{hyperref}
\hypersetup{
  colorlinks=true, linkcolor=black, citecolor=linkblue, urlcolor=linkblue,
  pdflang={th},
}

% Thai labels for auto-generated headings.
\renewcommand{\contentsname}{สารบัญ (Table of Contents)}
\renewcommand{\tablename}{ตาราง}
\renewcommand{\figurename}{รูปที่}

:
%    1) \hypersetup{pdftitle={...}}  — ชื่อเอกสารใน metadata ของ PDF
%    2) \begin{document}
%  (pdftitle ถูกย้ายออกจากไฟล์นี้ เพราะเป็นค่าเฉพาะของแต่ละเอกสาร)
%
%  แก้ที่นี่ที่เดียว = มีผลกับทุกเอกสาร ห้ามก๊อป preamble ไปวางซ้ำ
% =============================================================================
% =============================================================================
%  main.tex  —  เอกสารข้อเสนอโครงการ (Academic Project Proposal)
%  ระบบบริหารจัดการการยืม–คืนอุปกรณ์มหาวิทยาลัย (ULMs)
%
%  IMPORTANT: This document contains Thai text and MUST be compiled with
%  XeLaTeX (not pdflatex). Thai needs a Unicode font + ICU line-breaking.
%      Build:  latexmk main.tex        (uses .latexmkrc -> xelatex)
%      or:     xelatex main.tex        (run twice for the table of contents)
% =============================================================================

\documentclass[11pt, a4paper]{article}

% ---- Fonts & Thai support ----------------------------------------------------
% fontspec lets XeLaTeX use any system font. Laksaman is a complete serif face
% that covers BOTH Thai and Latin glyphs, so one \setmainfont handles the whole
% (Thai + English) document without font switching.
\usepackage{fontspec}

% Thai has no spaces between words, so TeX cannot find line-break points on its
% own. These two XeTeX primitives switch on ICU's dictionary-based Thai word
% segmentation, which is what makes Thai paragraphs wrap correctly.
\XeTeXlinebreaklocale "th"
\XeTeXlinebreakskip = 0pt plus 0.1pt

\setmainfont{Laksaman}

% ---- Page geometry & paragraph style ----------------------------------------
\usepackage[margin=2.5cm]{geometry}
\usepackage{parskip}          % block paragraphs (no indent, gap between) — common
                              % and more readable for Thai documents
\linespread{1.15}             % a little extra leading; Thai has tall tone marks

% ---- Colours (monochrome / grayscale palette) ------------------------------
\usepackage[table]{xcolor}    % 'table' option enables \rowcolor for table headers
% Monochrome (grayscale) palette. Names kept as-is so every reference still
% works; only the tones changed from the original green/blue accents.
\definecolor{kugreen}{HTML}{1A1A1A}   % near-black — section headings & title
\definecolor{kudark}{HTML}{333333}    % dark gray — subsection/subsubsection titles
\definecolor{headbg}{HTML}{EAEAEA}    % light gray background for table header rows
\definecolor{linkblue}{HTML}{404040}  % medium-dark gray — links / citations

% ---- Section heading styling -------------------------------------------------
% 'nobottomtitles*' stops a heading from being printed near the bottom of a
% page: if less than \bottomtitlespace of room is left, the heading (and its
% following text) moves to the next page. This is titlesec's built-in, robust
% replacement for the earlier \needspace hack — the latter injected vertical
% glue that clashed with longtable page-splitting ("Infinite glue shrinkage").
\usepackage[nobottomtitles*]{titlesec}
\titleformat{\section}
  {\color{kugreen}\Large\bfseries}{\thesection}{0.6em}{}
\titleformat{\subsection}
  {\color{kudark}\large\bfseries}{\thesubsection}{0.5em}{}
\titleformat{\subsubsection}
  {\color{kudark}\normalsize\bfseries}{\thesubsubsection}{0.5em}{}
\titlespacing*{\section}{0pt}{1.4\baselineskip}{0.5\baselineskip}

% Minimum room a heading needs at the bottom of a page; if only ~1–3 lines are
% left (less than this), the heading is bumped to the next page instead of being
% stranded above its table/text.
\renewcommand{\bottomtitlespace}{4\baselineskip}

% ---- Hierarchical indentation by heading level ------------------------------
% Each heading level shifts BOTH the heading and the body text under it further
% to the right: section = 0, subsection (x.y) = 1 step, subsubsection (x.y.z) =
% 2 steps. We do this by having the sectioning commands set \leftskip (a global
% left margin that carries through every following line until the next heading
% resets it). \par first, so the *previous* paragraph is not re-indented.
% \leftskip indents plain paragraphs; \@totalleftmargin makes list environments
% (itemize/enumerate) start from the same indent instead of the page margin.
\makeatletter
\newlength{\lvlindent}
\setlength{\lvlindent}{1.8em}          % one indent step
\let\ulmsSecIndent\section
\renewcommand{\section}{\par\global\leftskip=0pt\global\@totalleftmargin=0pt\relax\ulmsSecIndent}
\let\ulmsSubsecIndent\subsection
\renewcommand{\subsection}{\par\global\leftskip=\lvlindent\global\@totalleftmargin=\lvlindent\relax\ulmsSubsecIndent}
\let\ulmsSubsubsecIndent\subsubsection
\renewcommand{\subsubsection}{\par\global\leftskip=2\lvlindent\global\@totalleftmargin=2\lvlindent\relax\ulmsSubsubsecIndent}
\makeatother

% Indent the headings themselves to match (titlesec ignores \leftskip for the
% heading line, so its left margin is set here via the first \titlespacing arg).
\titlespacing*{\subsection}   {\lvlindent} {1.0\baselineskip}{0.4\baselineskip}
\titlespacing*{\subsubsection}{2\lvlindent}{0.8\baselineskip}{0.3\baselineskip}

% ---- Lists (tighter, custom bullets) ----------------------------------------
\usepackage{enumitem}
\setlist[itemize]{leftmargin=1.6em, itemsep=2pt, topsep=3pt}
\setlist[enumerate]{leftmargin=1.9em, itemsep=2pt, topsep=3pt}

% ---- Tables ------------------------------------------------------------------
\usepackage{array}
\usepackage{booktabs}         % \toprule / \midrule / \bottomrule
\usepackage{tabularx}         % auto-width columns that wrap long Thai text
\usepackage{longtable}        % tables that can break across pages
\usepackage{multirow}         % merged cells for the team-structure table

% Left-aligned wrapping column types:
%   L{width} — fixed-width ragged-right paragraph column
%   Y        — flexible ragged-right X column (for tabularx)
\newcolumntype{L}[1]{>{\raggedright\arraybackslash}p{#1}}
\newcolumntype{Y}{>{\raggedright\arraybackslash}X}
\renewcommand{\arraystretch}{1.3}

% Helpers for a styled header row inside tables.
%   \hrow  — starts a header row and shades it (must be FIRST in the row;
%            colortbl allows only one \rowcolor per row).
%   \thead — bold header-cell text.
\newcommand{\hrow}{\rowcolor{headbg}}
\newcommand{\thead}[1]{\textbf{#1}}

% \begin{keeptable} … \end{keeptable} — keeps a table and the bold label above
% it on the SAME page (used by the data dictionary in §5). A tabularx cannot
% split across pages, so LaTeX is free to break between the label and its
% table, stranding the label alone at the bottom of a page; wrapping the pair
% in a minipage makes them one unbreakable block.
%   * \leftskip is zeroed inside: the minipage box is already placed at the
%     current heading indent, so leaving it set would indent the content twice
%     and push the table past the right margin.
%   * \parskip is restored because minipage's \@parboxrestore zeroes it, which
%     would glue the table straight onto the label line.
% Tables inside must use \linewidth (= the minipage width), not \textwidth.
\newenvironment{keeptable}
  {\par\medskip\noindent
   \begin{minipage}{\dimexpr\textwidth-\leftskip\relax}%
   \leftskip=0pt\relax
   \setlength{\parskip}{0.5\baselineskip plus 2pt}}
  {\end{minipage}\par\medskip}

% ---- Table of contents styling ----------------------------------------------
% By default LaTeX gives dotted leaders to subsections but NOT to sections
% (they print bold with just a right-aligned page number). tocloft lets us turn
% on dots for the section level too, so every TOC line runs "….." to its page.
\usepackage{tocloft}
\renewcommand{\cftsecdotsep}{\cftdotsep}   % dotted leaders for section entries
\renewcommand{\cftsecleader}{\cftdotfill{\cftsecdotsep}}
\renewcommand{\cftsecpagefont}{\normalfont}  % keep page numbers un-bold

% ---- Figures -----------------------------------------------------------------
\usepackage{graphicx}
\graphicspath{{figures/}}

% Landscape (rotated) full-page floats — needed by the ER diagram in §5, which
% is far wider than it is tall and would be unreadable scaled to \textwidth.
\usepackage{rotating}

% ---- Flowcharts / diagrams (TikZ) -------------------------------------------
% Used for the borrow–return workflow diagram (§5). Libraries: geometric shapes
% for decision diamonds, arrows.meta for arrowheads, positioning for relative
% node placement, calc/fit/backgrounds for routing loops and grouping.
\usepackage{tikz}
\usetikzlibrary{shapes.geometric, arrows.meta, positioning, calc, fit, backgrounds}

% ---- Links & bookmarks -------------------------------------------------------
\usepackage{hyperref}
\hypersetup{
  colorlinks=true, linkcolor=black, citecolor=linkblue, urlcolor=linkblue,
  pdflang={th},
}

% Thai labels for auto-generated headings.
\renewcommand{\contentsname}{สารบัญ (Table of Contents)}
\renewcommand{\tablename}{ตาราง}
\renewcommand{\figurename}{รูปที่}



\hypersetup{pdftitle={รายงานความก้าวหน้าโครงการ ULMs ครั้งที่ 4}}

% ---- ย่อหน้า (เหมือนฉบับที่ 3) ------------------------------------------------
\setlength{\parindent}{1.5em}
\setlength{\parskip}{0.35\baselineskip plus 2pt}
\titlespacing{\section} {0pt} {1.4\baselineskip}{0.5\baselineskip}
\titlespacing{\subsection} {\lvlindent} {1.0\baselineskip}{0.4\baselineskip}
\titlespacing{\subsubsection}{2\lvlindent}{0.8\baselineskip}{0.3\baselineskip}

\setlist[itemize]{leftmargin=1.6em, itemsep=2pt, topsep=0.55\baselineskip}
\setlist[enumerate]{leftmargin=1.9em, itemsep=2pt, topsep=0.55\baselineskip}

\setlength{\emergencystretch}{3em}
\newcommand{\zb}{\hspace{0pt}}

\usepackage{float}

% ---- ตัวแปรของรอบรายงาน ------------------------------------------------------
\newcommand{\reportdate}{17 กันยายน 2569}
\newcommand{\reportperiod}{11 – 17 กันยายน 2569}
\newcommand{\actualpct}{79\%}
\newcommand{\planpct}{86\%}
\newcommand{\prevactualpct}{75\%}

\begin{document}

% =============================================================================
% TITLE PAGE
% =============================================================================
\begin{titlepage}
\centering
\vspace*{2.8cm}

{\Large\bfseries\color{kugreen} รายงานความก้าวหน้าโครงการ ครั้งที่ 4\par}
\vspace{0.6cm}
{\LARGE\bfseries\color{kugreen} ระบบบริหารจัดการการยืม–คืนอุปกรณ์มหาวิทยาลัย\par}
\vspace{0.35cm}
{\Large\color{kudark} University Lending Management System (ULMs)\par}

\vspace{1.8cm}
\rule{0.72\textwidth}{0.8pt}
\vspace{0.9cm}

\begin{tabular}{@{}r@{\hspace{1.2em}}l@{}}
\textbf{รอบรายงาน} & \reportperiod \\[0.4em]
\textbf{วันที่รายงาน} & \reportdate \\[0.4em]
\textbf{ขอบเขตการนับ} & งานพัฒนาระบบ Sprint 0–5 (ไม่รวมงานก่อน 24 ก.ค.) \\[0.4em]
\textbf{สถานะ Sprint} & วันที่ 7 จาก 14 ของ Sprint 4 (Sprint ที่ 5 จาก 6) \\[0.4em]
\textbf{ความก้าวหน้า} & \actualpct\ เทียบแผน \planpct\ \ (ต่ำกว่าแผน 7\%) \\[0.4em]
\textbf{งบประมาณ} & ใช้ไป 0 บาท จากวงเงินประเมิน 5{,}000 บาท \\[0.4em]
\textbf{กำหนดเสร็จ} & 1 ตุลาคม 2569 (ยังไม่มีการเลื่อน) \\
\end{tabular}

\vspace{0.9cm}
\rule{0.72\textwidth}{0.8pt}

\vfill
{\large\color{kudark} คณะทำงาน ``miro ปั่น 14 บาท''\par}
\vspace{0.3cm}
{\color{kudark} ภาควิชาวิศวกรรมคอมพิวเตอร์ คณะวิศวกรรมศาสตร์\par}
{\color{kudark} มหาวิทยาลัยเกษตรศาสตร์\par}
\vspace{1.2cm}
\end{titlepage}

% =============================================================================
\section{บทสรุป}

รอบรายงานนี้เป็นครึ่งแรกของ Sprint 4 พอดี คือผ่านไป 7 วันจาก 14 วัน ความก้าวหน้าเพิ่มจาก \prevactualpct\ เป็น \actualpct\ ขณะที่แผนขยับจาก 78\% เป็น \planpct\ ระยะห่างจากแผนจึงกว้างขึ้นจาก 3\% เป็น 7\%

เรื่องสำคัญที่สุดของรอบนี้คืองานต่อเชื่อมที่ค้างมาสองรอบเกิดขึ้นจริงแล้ว ปุ่มส่งคำขอของผู้ยืมเรียก loan.create และเขียนลงฐานข้อมูลได้ และปุ่มยกเลิกคำขอเรียก loan.cancel ทั้งสองรายการคือสาเหตุหลักที่ทำให้จุดวัดของรอบก่อนไม่ผ่าน คณะทำงานได้ตรวจสอบด้วยการรันระบบจริงบนฐานข้อมูลที่ seed ใหม่ และชุดทดสอบผ่านเบราว์เซอร์ยืนยันการส่งคำขออุปกรณ์สองชิ้นผ่าน loan.create สำเร็จ

หน้ารับของเป็นข้อยกเว้นที่ต้องแจ้ง หน้าจอถูกต่อเข้ากับ backend แล้วก็จริง แต่เรียกคำสั่งสามตัวที่ฝั่งเซิร์ฟเวอร์ไม่มี คือ loan.requestPickupImageUpload loan.attachPickupImage และ loan.finalizePickup ส่วนคำสั่งที่เซิร์ฟเวอร์มีจริงคือ image.requestUpload ซึ่งเป็นคนละชื่อ หน้าจอจึงผ่านการตรวจชนิดข้อมูลแต่ใช้งานจริงไม่ได้ สาเหตุคือหน้าเว็บอ้างอิงไฟล์สัญญาที่เขียนด้วยมือ ซึ่งประกาศคำสั่งเหล่านั้นไว้เองโดยไม่มีอะไรตรวจว่าเซิร์ฟเวอร์มีจริงหรือไม่ ตรงกับความเสี่ยงที่จุดวัดข้อ 3 พยายามปิด

อย่างไรก็ตาม จุดวัดสี่ข้อที่รายงานครั้งที่ 3 ประกาศไว้ยังไม่ผ่านสามข้อ เพราะแต่ละข้อกำหนดขอบเขตไว้กว้างกว่าที่ทำได้จริง ข้อแรกขอทั้งปุ่มส่งคำขอและหน้าส่งมอบของเจ้าหน้าที่ แต่หน้าส่งมอบยังเป็นหน้าว่างอยู่ ข้อที่สองขอให้วงจรยืม–คืนเดินครบเส้นทางพร้อมชุดทดสอบในระบบตรวจอัตโนมัติ แต่ชุดทดสอบที่มีครอบคลุมเฉพาะขั้นส่งคำขอ และระบบตรวจอัตโนมัติยังไม่ได้เรียกชุดทดสอบผ่านเบราว์เซอร์เลย ข้อที่สามขอให้หน้าเว็บใช้ไฟล์สัญญาที่สร้างจากโค้ด แต่ไฟล์ที่ใช้อยู่ยังเขียนด้วยมือ และหัวไฟล์ระบุเองว่าเป็นของชั่วคราว

งานทุกข้อที่ผ่านในรอบนี้ช้ากว่ากำหนดที่ตั้งไว้ 2 ถึง 4 วัน ปุ่มส่งคำขอกำหนด 12 กันยายน แต่รวมเข้า main วันที่ 14 และ 16 กันยายน ส่วนชุดทดสอบผ่านเบราว์เซอร์ของการยืมกำหนด 14 กันยายน แต่รวมเข้าวันที่ 17 กันยายน ซึ่งเป็นวันรายงานนี้

มิติที่ต้องเฝ้าระวังยังเป็น Test เหมือนสองรอบก่อน แต่รอบนี้ขยับมากที่สุดในบรรดาทุกมิติ คือ 12 จุด จากชุดทดสอบฝั่ง backend ที่เพิ่มจาก 234 เป็น 266 กรณี และชุดทดสอบผ่านเบราว์เซอร์ที่เพิ่มจาก 11 เป็น 13 กรณี มิติที่แทบไม่ขยับกลับกลายเป็น Back-End Dev เพราะจำนวน procedure บน main เท่าเดิมทุกตัว งานฝั่ง backend ในรอบนี้เป็นการแก้ความถูกต้องและเพิ่มชุดทดสอบ ไม่ใช่การเพิ่มความสามารถใหม่

% =============================================================================
\section{แผนภาพรวมทั้ง 6 Sprint}
\label{sec:overview}

\noindent\begin{tabularx}{\dimexpr\textwidth-\leftskip\relax}{@{} L{2.9cm} Y L{2.2cm} @{}}
\toprule
\hrow \thead{Sprint} & \thead{ธีมและงานหลัก} & \thead{สถานะ} \\
\midrule
\textbf{Sprint 0} \newline 24 – 30 ก.ค. & วางฐาน — ปิดขอบเขต · ติดตั้งโครงสร้างพื้นฐานทั้งสองฝั่ง · database schema แกนกลาง · ระบบออกแบบและ mock & ผ่าน \\
\textbf{Sprint 1} \newline 31 ก.ค. – 13 ส.ค. & ล็อก contract และเปิดทาง — ล็อก contract API · ระบบเข้าสู่ระบบและออกจากระบบ · CI · เกณฑ์การยอมรับของวงจรหลัก & ผ่าน \\
\textbf{Sprint 2} \newline 14 – 27 ส.ค. & วงจรยืม–คืนของอุปกรณ์เคลื่อนย้ายได้ — รายการอุปกรณ์และจำนวนคงเหลือ · สร้างคำขอ อนุมัติ ส่งมอบ รับคืน · ข้อมูล seed · เลิกใช้ mock ในเส้นทางหลัก & ไม่ผ่านเกณฑ์ \newline (ยกงานไป Sprint 3) \\
\textbf{Sprint 3} \newline 28 ส.ค. – 10 ก.ย. & กฎธุรกิจรายระดับและการจองสถานที่ (T3) — อนุมัติตามระดับอุปกรณ์ · ผูกหมายเลขชิ้น · สล็อตเวลาและการจองห้อง · ตรวจสภาพและระบบเครดิต & ไม่ผ่านเกณฑ์ \newline (ยกงาน 4 วันไป Sprint 4) \\
\textbf{Sprint 4} \newline 11 – 24 ก.ย. & ข้ามหน่วยงานและงานขัดเงา — ตะกร้าข้ามหน่วยงาน · จองล่วงหน้า · ต่ออายุ · หน้าจอที่เหลือ · การค้นหาและการรองรับมือถือ \newline สี่วันแรกรับงานค้างของ Sprint 3 ก่อน & กำลังทำ \newline (ผ่านไป 7 จาก 14 วัน งานค้างที่ยกมาทำได้บางส่วน) \\
\textbf{Sprint 5} \newline 25 ก.ย. – 1 ต.ค. & หยุดเพิ่มความสามารถ — ทดสอบถดถอย · นำขึ้นใช้งานจริง · คู่มือผู้ใช้และคู่มือติดตั้ง · ซ้อมนำเสนอ & ตามแผน \\
\bottomrule
\end{tabularx}

% =============================================================================
\section{สถานะโครงการโดยรวม}

\noindent\begin{tabularx}{\dimexpr\textwidth-\leftskip\relax}{@{} L{1.5cm} L{3.0cm} L{2.4cm} Y @{}}
\toprule
\hrow \thead{ด้าน} & \thead{แผน} & \thead{ผลจริง} & \thead{การประเมิน} \\
\midrule
เวลา & ผ่านไป 7.9 จาก 10 สัปดาห์ (79\%) & ผ่านไป 7.9 จาก 10 สัปดาห์ (79\%) & ตามปฏิทิน วันเสร็จยังเป็น 1 ตุลาคม 2569 เหลือเวลาอีก 2 สัปดาห์ ซึ่ง Sprint 5 ทั้งช่วงสงวนไว้สำหรับการทดสอบและนำขึ้นใช้งาน ไม่ใช่การเพิ่มความสามารถ \\
งาน & \planpct & \actualpct & ต่ำกว่าแผน 7\% กว้างขึ้นจาก 3\% ในรอบก่อน งานต่อเชื่อมฝั่งผู้ยืมเกิดขึ้นแล้วเป็นครั้งแรก แต่ความสามารถใหม่ฝั่ง backend หยุดนิ่งทั้งรอบ รายละเอียดอยู่ในหัวข้อ \ref{sec:dimension} \\
เงิน & 5{,}000 บาท & 0 บาท (0\%) & ยังไม่มีรายจ่าย รายการที่ประเมินไว้คือค่าเช่า Cloud/Hosting และฐานข้อมูลกลาง ซึ่งจะเกิดขึ้นเมื่อนำระบบขึ้นใช้งานจริงใน Sprint 5 \\
\bottomrule
\end{tabularx}

% =============================================================================
\clearpage
\section{ความก้าวหน้าแยกตามมิติของงาน}
\label{sec:dimension}

ค่าในคอลัมน์ ``จริง'' นับเฉพาะงานที่รวมเข้า main แล้ว วันรายงานอยู่กึ่งกลาง Sprint 4 พอดี ค่าในคอลัมน์ ``แผน'' จึงสอดแทรกครึ่งทางระหว่างแถว S3 กับ S4 ของภาคผนวก

\noindent\begin{tabularx}{\dimexpr\textwidth-\leftskip\relax}{@{} L{3.2cm} L{1.1cm} L{1.0cm} L{1.0cm} L{1.2cm} Y @{}}
\toprule
\hrow \thead{มิติงาน} & \thead{น้ำหนัก} & \thead{แผน} & \thead{จริง} & \thead{ผลต่าง} & \thead{ข้อเท็จจริงประกอบ} \\
\midrule
Back-End Design \newline (API + Database Model) & 10\% & 100\% & 97\% & $-3$ & contract คงที่ 10 domain 86 procedure ตลอดทั้งรอบ คอลัมน์เวลาเปลี่ยนเป็นชนิด timestamptz ที่เก็บเขตเวลาไว้ด้วย ยังเหลือการออกแบบการอุทธรณ์และการจองสล็อตเวลา \\
Front-End Design \newline (UI/UX + Business Process) & 10\% & 96\% & 96\% & $0$ & หน้าจอที่มีเนื้อหาจริง 27 จาก 29 หน้า หน้าห้องปรับเป็นช่องละ 30 นาที จัดกลุ่มตามอาคาร เหลือหน้าว่าง 2 หน้า คือส่งมอบของและส่งออกรายงาน \\
Back-End Dev & 30\% & 88\% & 79\% & $-9$ & จำนวน procedure บน main เท่ารอบก่อนทุกตัว งานรอบนี้เป็นการแก้ความถูกต้องของเวลา และกันการลดบทบาทจนหน่วยงานไม่เหลือผู้ดูแล ยังเหลือการอุทธรณ์ซึ่งไม่มีโค้ดเลย และ cron job อีก 3 งาน \\
Front-End Dev & 25\% & 88\% & 80\% & $-8$ & ปุ่มส่งคำขอและปุ่มยกเลิก เรียก backend ได้แล้ว ซึ่งเป็นงานที่ค้างมาสองรอบ ส่วนหน้ารับของเรียกคำสั่งที่เซิร์ฟเวอร์ไม่มี จึงยังใช้งานไม่ได้ ยังเหลือปุ่มขอต่ออายุ หน้าห้อง 3 หน้าและหน้าอุทธรณ์ 2 หน้า \\
Test (Front-End + Back-End) & 15\% & 66\% & 47\% & $-19$ & ชุดทดสอบฝั่ง backend เพิ่มจาก 234 เป็น 266 กรณี ชุดผ่านเบราว์เซอร์เพิ่มจาก 11 เป็น 13 และมีชุดของการส่งคำขอยืมเป็นครั้งแรก service ที่มีชุดทดสอบเป็น 10 จาก 23 แต่ระบบตรวจอัตโนมัติยังไม่เรียกชุดผ่านเบราว์เซอร์ \\
Documentation & 10\% & 82\% & 85\% & $+3$ & SDS ปรับเป็นฉบับ 1.2 นับใหม่จากโค้ดของวันรายงาน พร้อมแก้เรื่องการระงับสิทธิ์การยืมให้ตรงแผน และเพิ่มรายงานผลตรวจความสอดคล้องทั้งเว็บ ซึ่งข้อเสนอแก้ SRS ในรายงานนั้นยังรอการตัดสิน \\
\midrule
\hrow \thead{รวมถ่วงน้ำหนัก} & \thead{100\%} & \thead{\planpct} & \thead{\actualpct} & \thead{$-7$} & \\
\bottomrule
\end{tabularx}

มิติที่ต้องเฝ้าระวังยังเป็น Test ซึ่งห่างจากแผนมากขึ้นอีกจาก $-17$ เป็น $-19$ แต่สาเหตุเปลี่ยนไปจากสองรอบก่อน รอบนี้ Test ขยับ 12 จุด มากที่สุดในบรรดาทุกมิติ ที่ยังห่างเพราะค่าแผนของ Sprint 4 ขยับเร็วกว่า คือจาก 52\% ไป 80\% ภายใน Sprint เดียว สิ่งที่ยังขาดคือชุดทดสอบของวงจรยืม–คืนที่เดินครบเส้นทาง ปัจจุบันครอบคลุมเฉพาะขั้นส่งคำขอ และการเรียกชุดทดสอบผ่านเบราว์เซอร์จากระบบตรวจอัตโนมัติ ซึ่งไฟล์ตั้งค่ามีอยู่ในรีโปแล้วแต่ยังไม่มีขั้นตอนใดเรียกใช้

มิติที่ถดถอยเชิงเปรียบเทียบคือ Back-End Dev ซึ่งรอบก่อนห่างจากแผนเพียง 1 จุด รอบนี้ห่าง 9 จุด เพราะค่าแผนขยับจาก 78\% ไป 88\% ขณะที่จำนวน procedure บน main เท่าเดิมทุกตัว คณะทำงานตรวจซ้ำด้วยการนับ decorator ทั้งที่ c8ed505 และที่ 0cb21a7 ได้ 86 เท่ากันทั้งสองจุด งานฝั่ง backend ที่เกิดขึ้นจริงจึงเป็นการแก้ความถูกต้องและการเพิ่มชุดทดสอบ ซึ่งนับเข้ามิติ Test ไม่ใช่ Back-End Dev

\begin{figure}[H]
\centering
\begin{tikzpicture}[x=0.062cm, y=1cm]
 \foreach \x in {0,20,40,60,80,100} {
 \draw[gray!25] (\x,-0.30) -- (\x,6.15);
 \node[below, font=\scriptsize] at (\x,-0.30) {\x};
 }
 \node[below, font=\scriptsize] at (50,-0.68) {\% ความก้าวหน้าของมิตินั้น};
 \draw[thick] (0,-0.30) -- (0,6.15);

 \foreach \y/\name/\wt/\plan/\act in {%
 5.6/{Back-End Design}/10/100/97,
 4.6/{Front-End Design}/10/96/96,
 3.6/{Back-End Dev}/30/88/79,
 2.6/{Front-End Dev}/25/88/80,
 1.6/{Test}/15/66/47,
 0.6/{Documentation}/10/82/85} {
 \node[left, xshift=-0.16cm, font=\scriptsize, align=right] at (0,\y)
 {\name\\[-3pt]{\tiny น้ำหนัก \wt\%}};
 \fill[gray!45] (0,{\y+0.05}) rectangle (\plan,{\y+0.30});
 \fill[kugreen] (0,{\y-0.30}) rectangle (\act,{\y-0.05});
 \node[right, font=\tiny, gray!45!black] at (\plan,{\y+0.175}) {\plan};
 \node[right, font=\tiny, kugreen] at (\act,{\y-0.175}) {\act};
 }

 \fill[gray!45] (0,-1.36) rectangle (7,-1.54);
 \node[right, font=\scriptsize] at (9,-1.45) {แผน};
 \fill[kugreen] (0,-1.76) rectangle (7,-1.94);
 \node[right, font=\scriptsize] at (9,-1.85) {ผลจริง};
\end{tikzpicture}
\vspace{0.6em}
\caption{ความก้าวหน้ารายมิติ ณ \reportdate\ ซึ่งเป็นวันที่ 7 จาก 14 ของ Sprint 4
 เทียบกับค่าเป้าหมายตามแผนที่สอดแทรกครึ่งทางระหว่างขอบ Sprint 3 กับ Sprint 4
 ตัวเลขใต้ชื่อมิติคือน้ำหนักที่ใช้ถ่วงเมื่อรวมเป็นค่าเดียว
 แท่งผลจริงนับเฉพาะงานที่รวมเข้า main แล้ว
 มิติที่ห่างจากแผนมากที่สุดคือ Test ที่ $-19$ รองลงมาคือ Back-End Dev ที่ $-9$
 ส่วน Documentation เดินนำแผนอยู่เล็กน้อยที่ $+3$}
\label{fig:dims}
\end{figure}

% =============================================================================
\clearpage
\section{กราฟความก้าวหน้าเทียบแผน}

เส้นแผนสร้างจากแผน Sprint โดยตรง คือกำหนดค่าเป้าหมายรายมิติ ณ วันสิ้นสุดของแต่ละ Sprint แล้วคูณน้ำหนักรวมกันเป็นค่าสะสม ตารางค่าทั้งหมดอยู่ในภาคผนวก แกนนอนเริ่มที่ 24 กรกฎาคม 2569 ซึ่งเป็นวันเริ่มลงมือพัฒนา

ความชันของเส้นผลจริงในรอบนี้คือ 4 จุดในเจ็ดวัน ลดลงจากรอบก่อนที่ 5 จุดในหกวัน ขณะที่เส้นแผนในช่วง Sprint 4 ชันกว่าทุกช่วงที่ผ่านมา คือ 16 จุดใน 14 วัน ระยะห่างจึงกว้างขึ้นเร็วกว่าทุกรอบที่ผ่านมา

\begin{figure}[H]
\centering
\begin{tikzpicture}[x=1.20cm, y=0.058cm]
 \foreach \y in {0,20,40,60,80,100}
 \draw[gray!22] (0,\y) -- (10.2,\y);

 \foreach \x in {0.86, 2.86, 4.86, 6.86, 8.86}
 \draw[gray!40, dashed, thin] (\x,0) -- (\x,108);

 \foreach \x/\lbl in {0.43/{S0}, 1.86/{S1}, 3.86/{S2}, 5.86/{S3}, 7.86/{S4}, 9.36/{S5}}
 \node[font=\scriptsize\bfseries, gray!50!black] at (\x,114) {\lbl};

 \draw[->, thick] (0,0) -- (10.6,0);
 \draw[->, thick] (0,0) -- (0,122) node[above, font=\scriptsize] {\% ความก้าวหน้าสะสม};
 \foreach \y in {0,20,40,60,80,100}
 \draw (0,\y) -- (-0.10,\y) node[left, font=\scriptsize] {\y};

 \foreach \x/\lbl in {0/{24 ก.ค.}, 0.86/{30 ก.ค.}, 2.86/{13 ส.ค.},
 4.86/{27 ส.ค.}, 6.86/{10 ก.ย.}, 8.86/{24 ก.ย.}, 9.86/{1 ต.ค.}}
 {\draw (\x,0) -- (\x,-3);
 \node[rotate=45, anchor=north east, font=\scriptsize] at (\x,-4) {\lbl};}

 % ---- เส้นแผน ----
 \draw[thick, kudark, densely dashed]
 plot coordinates {(0,1) (0.86,12) (2.86,42) (4.86,65) (6.86,78) (8.86,94) (9.86,100)};
 \foreach \p in {(0,1), (0.86,12), (2.86,42), (4.86,65), (6.86,78), (8.86,94), (9.86,100)}
 \fill[kudark] \p circle (1.6pt);

 % ---- เส้นผลจริง — จุดคือวันที่วัด 7.86 = 17 ก.ย. = วันรายงานนี้ ----
 \draw[very thick, kugreen]
 plot coordinates {(0,0) (2.0,24) (2.43,32) (2.86,43) (4.29,43) (4.86,52) (6.0,70) (6.86,75) (7.86,79)};
 \foreach \p in {(0,0), (2.0,24), (2.43,32), (2.86,43), (4.29,43), (4.86,52), (6.0,70), (6.86,75), (7.86,79)}
 \fill[kugreen] \p circle (2.2pt);

 \draw[gray!70, thick, |-|] (7.86,79) -- (7.86,86);
 \node[right, font=\tiny, gray!70!black] at (7.92,82.5) {$-7$};

 \node[anchor=south east, font=\scriptsize, gray!55!black, align=right,
 fill=white, inner sep=1.5pt] at (6.80,77)
 {10 ก.ย.\\75\% · แผน 78\%};
 \node[anchor=north west, font=\scriptsize, kugreen, align=left,
 fill=white, inner sep=1.5pt] at (7.95,75)
 {17 ก.ย. (วันรายงาน)\\ผลจริง 79\% · แผน 86\%};

 \fill[white] (0.55,72) rectangle (4.60,96);
 \draw[thick, kudark, densely dashed] (0.70,89) -- (1.40,89);
 \node[right, font=\scriptsize] at (1.50,89) {แผน (Planned)};
 \draw[very thick, kugreen] (0.70,79) -- (1.40,79);
 \node[right, font=\scriptsize] at (1.50,79) {ผลจริง (Actual)};
\end{tikzpicture}
\caption{กราฟ เปรียบเทียบความก้าวหน้าตามแผนกับผลการดำเนินงานจริง
 แกนนอนวางตามขอบ Sprint ทั้ง 6 ช่วง (เส้นประแนวตั้ง)
 จุดบนเส้นผลจริงคือวันที่วัดความก้าวหน้า
 วันรายงาน 17 กันยายน 2569 อยู่กึ่งกลาง Sprint 4 ไม่ตรงขอบ Sprint
 ค่าแผน ณ วันนั้นจึงอ่านจากเส้นแผนที่ลากระหว่างขอบ Sprint 3 กับ Sprint 4
 ช่วงราบระหว่าง 13 ถึง 23 สิงหาคม คือช่วงสอบกลางภาค}
\label{fig:scurve}
\end{figure}

% =============================================================================
\clearpage
\section{ผลของจุดวัดที่รายงานครั้งที่ 3 ตั้งไว้}
\label{sec:review}

รายงานครั้งที่ 3 ประกาศจุดวัดไว้สี่ข้อ และระบุว่าถ้าข้อ 1 หรือข้อ 2 ไม่ผ่าน จะตัดตะกร้าข้ามหน่วยงานและการจองล่วงหน้าออกจาก Sprint 4 ก่อน เพื่อรักษาเวลาทดสอบของ Sprint 5 ตารางนี้คือผลของทั้งสี่ข้อ

\noindent\begin{tabularx}{\dimexpr\textwidth-\leftskip\relax}{@{} Y L{1.7cm} L{5.4cm} L{1.9cm} @{}}
\toprule
\hrow \thead{จุดวัด} & \thead{กำหนด} & \thead{สภาพจริง ณ วันรายงาน} & \thead{ผล} \\
\midrule
ผู้ยืมส่งคำขอผ่านหน้าจอ ลงฐานข้อมูลจริง และหน้าส่งมอบใช้งานได้ & 12 ก.ย. & ส่งคำขอทำได้แล้ว รวมเข้า main วันที่ 14 และ 16 ก.ย. ช้ากว่ากำหนด 2 ถึง 4 วัน แต่หน้าส่งมอบของเจ้าหน้าที่ยังเป็นหน้าว่าง & ไม่ผ่าน \newline (ผ่านครึ่งเดียว) \\
วงจรยืม–คืน เดินผ่านหน้าจอจริง และมีชุดทดสอบ ในระบบตรวจอัตโนมัติ & 14 ก.ย. & ชุดทดสอบผ่านเบราว์เซอร์ของการยืมรวมเข้า main วันที่ 17 ก.ย. แต่ครอบคลุมเฉพาะขั้นส่งคำขอ ยังไม่มีขั้นอนุมัติ ส่งมอบ รับคืน และตรวจสภาพ และระบบตรวจอัตโนมัติยังไม่เรียกชุดนี้ & ไม่ผ่าน \\
ฝั่ง Front-End ใช้ไฟล์สัญญา ที่สร้างจากโค้ด & 14 ก.ย. & หน้าเว็บรวมสัญญามาไว้ที่ไฟล์เดียวแล้ว ซึ่งเป็นความคืบหน้า แต่ไฟล์นั้นยังเขียนด้วยมือ และหัวไฟล์ระบุเองว่าเป็นของชั่วคราวที่ต้องแทนที่ด้วยของที่ backend สร้าง & ไม่ผ่าน \\
ไม่เหลือข้อมูลจำลอง และหน้าว่างในระบบ & 24 ก.ย. & เหลือหน้าว่าง 2 หน้า และเหลือ 5 หน้าที่อยู่ในโดเมนซึ่งยังไม่มี backend รองรับเลย คือห้อง 3 หน้าและอุทธรณ์ 2 หน้า เท่ากับรอบก่อน & ยังไม่ถึงกำหนด \\
\bottomrule
\end{tabularx}

ผลของจุดวัดในรอบนี้เหมือนรอบก่อนในเชิงตัวเลข คือไม่ผ่านสามข้อและยังไม่ถึงกำหนดหนึ่งข้อ แต่สาเหตุต่างกันโดยสิ้นเชิง รอบก่อนไม่ผ่านเพราะงานไม่ได้เริ่ม รอบนี้ไม่ผ่านเพราะงานเริ่มแล้วแต่ทำได้ไม่ครบขอบเขตที่จุดวัดกำหนด และทุกข้อที่ทำได้ช้ากว่ากำหนด

ข้อสังเกตที่คณะทำงานบันทึกไว้คือ จุดวัดสองข้อแรกรวมงานสองชิ้นไว้ในข้อเดียว ทำให้ผลออกมาเป็น ``ไม่ผ่าน'' ทั้งที่ครึ่งหนึ่งของแต่ละข้อทำได้แล้ว จุดวัดของรอบถัดไปจึงแยกเป็นข้อละหนึ่งงาน เพื่อให้ผลสะท้อนความคืบหน้าได้ตรงกว่า

เนื่องจากข้อ 1 และข้อ 2 ไม่ผ่าน คณะทำงานจึงเรียกใช้แผนสำรองที่ประกาศไว้ คือตัดตะกร้าข้ามหน่วยงานและการจองล่วงหน้าออกจาก Sprint 4 รายละเอียดอยู่ในหัวข้อ \ref{sec:next}

% =============================================================================
\clearpage
\section{ผลงานที่ส่งมอบในรอบนี้}

รอบนี้มี 29 commit รวมเข้า main ผ่าน 10 pull request จากผู้ร่วมงาน 5 คน

\subsection{ฝั่ง frontend}

\begin{itemize}
\item ปุ่มส่งคำขอของผู้ยืมเรียก loan.create เขียนลงฐานข้อมูลจริง และรายงานผลรายชิ้น
\item ปุ่มยกเลิกคำขอเรียก loan.cancel พร้อมกล่องยืนยันก่อนยกเลิก
\item หน้ารับของถูกต่อเข้ากับ backend แล้ว แต่เรียกคำสั่งที่เซิร์ฟเวอร์ไม่มี จึงยังใช้งานจริงไม่ได้
\item หน้ารายการคำขอของผู้ยืมอ่านจาก backend ทั้งหมด ข้อถดถอยที่รายงานครั้งที่ 3 แจ้งไว้ว่าคำขอที่เพิ่งส่งหายจากรายการ ไม่เกิดขึ้นอีก
\item หน้าห้องปรับตามมติทีม คือช่องเวลาละ 30 นาทีเริ่ม 07:00 และจัดกลุ่มตามอาคารแทนชนิดห้อง
\item หน้ารายการอุปกรณ์แก้การกรองให้ตรงหน่วยงานเจ้าของ และหน้ารายละเอียดตัดแถวที่ซ้ำออก
\item รวมสัญญาระหว่างสองฝั่งมาไว้ที่ไฟล์เดียว client จึงมีชนิดข้อมูลครบทั้งระบบ แม้ไฟล์นั้นยังเขียนด้วยมือ
\end{itemize}

\subsection{ฝั่ง backend}

\begin{itemize}
\item คอลัมน์เวลาทุกคอลัมน์ในฐานข้อมูลเปลี่ยนเป็นชนิด timestamptz ซึ่งเก็บเขตเวลาไว้กับค่า ค่าเวลาจากเครื่องคนละเขตเวลาจึงเทียบกันได้ตรง และกำหนดคืนไม่เลื่อนตามผู้ใช้
\item เพิ่มการป้องกันการเปลี่ยนบทบาทหรือปิดบัญชีจนหน่วยงานไม่เหลือผู้ดูแล ซึ่งเดิมระบบไม่เตือน
\item เพิ่มชุดทดสอบของการจัดการผู้ใช้และการจัดการอุปกรณ์ พร้อมไฟล์ตั้งค่าแยกสำหรับรันเฉพาะกลุ่ม
\item จำนวน procedure บน main ไม่เปลี่ยน ยังเป็น 86 procedure ใน 10 domain
\end{itemize}

\subsection{เครื่องมือและกระบวนการ}

\begin{itemize}
\item มีชุดทดสอบผ่านเบราว์เซอร์ของการส่งคำขอยืมเป็นครั้งแรก รวมเป็น 13 กรณีใน 4 ไฟล์
\item คณะทำงานจัดทำรายการความพร้อมของการทดสอบรายโมดูลไว้ในรีโป ครอบคลุมทั้ง 12 โมดูล แยกสถานะทั้งสองฝั่งและการต่อเชื่อม
\item ล้างสำเนารีโปที่ถูก commit ซ้อนเข้ามาโดยไม่ตั้งใจ พร้อมเพิ่มกติกาไม่ให้เกิดซ้ำ
\end{itemize}

\subsection{เอกสารและการออกแบบ}

\begin{itemize}
\item SDS ปรับเป็นฉบับ 1.2 นับสถานะใหม่จากโค้ดของวันรายงาน
\item แก้ข้อที่ฉบับ 1.1 เขียนตามโค้ดแทนที่จะเขียนตามแผน คือการระงับสิทธิ์การยืม ซึ่งแผนกำหนดให้บทลงโทษมีผลผ่านระดับเครดิตเท่านั้น ฟังก์ชันระงับสิทธิ์ในโค้ดถูกกำกับว่าอยู่นอกแผน
\item จัดทำรายงานผลตรวจความสอดคล้องทั้งเว็บ พร้อมฉบับย่อสำหรับให้ทีมอ่านเร็ว
\end{itemize}

% =============================================================================
\clearpage
\section{หลักฐานการทดสอบที่รันจริงในรอบนี้}
\label{sec:evidence}

ตัวเลขในตารางนี้ได้จากการรันบนเครื่องของคณะทำงานในวันรายงาน โดยดึงโค้ดจาก main ที่ 0cb21a7 สร้างฐานข้อมูลใหม่แล้ว seed ข้อมูลตั้งต้น จากนั้นสตาร์ต backend และ frontend เองก่อนรันชุดทดสอบผ่านเบราว์เซอร์ ไม่ได้อ่านจากผลของระบบตรวจอัตโนมัติ

\noindent\begin{tabularx}{\dimexpr\textwidth-\leftskip\relax}{@{} L{4.6cm} L{2.6cm} Y @{}}
\toprule
\hrow \thead{ชุดที่รัน} & \thead{ผล} & \thead{หมายเหตุ} \\
\midrule
ชุดทดสอบฝั่ง backend & ผ่าน 266 จาก 266 & 24 ชุดทดสอบ รายงานครั้งที่ 3 แจ้ง 234 กรณี ส่วนที่เพิ่มมาจากชุดจัดการผู้ใช้ ชุดจัดการอุปกรณ์ และชุดตรวจรูปแบบวันเวลา \\
ตรวจชนิดข้อมูลฝั่ง frontend & ผ่าน & ไม่มีข้อผิดพลาด \\
ชุดทดสอบส่วนประกอบฝั่ง frontend & ผ่าน 33 จาก 33 & 3 ไฟล์ ครอบคลุมหน้าจัดการผู้ใช้ รายการอุปกรณ์ และคลังอุปกรณ์ \\
ชุดตรวจเกณฑ์การยอมรับ & ผ่าน 49 จาก 49 & รายงานครั้งที่ 3 แจ้ง 44 ข้อ \\
ตรวจรูปแบบโค้ดฝั่ง frontend & 0 ข้อผิดพลาด 7 คำเตือน & คำเตือนทั้งหมดเป็นเรื่องการแยกไฟล์ของส่วนประกอบ ไม่กระทบการทำงาน \\
ชุดทดสอบผ่านเบราว์เซอร์ & ผ่าน 12 ข้าม 1 จาก 13 & กรณีที่ข้ามคือการเปิดหน้ารายละเอียดอุปกรณ์ ซึ่งข้ามตัวเองเมื่อหาปุ่มเปิดไม่พบ จึงไม่ยืนยันว่าเส้นทางนั้นใช้ได้ \\
\bottomrule
\end{tabularx}

ข้อควรระวังสองข้อที่คณะทำงานบันทึกไว้ ข้อแรก รายการความพร้อมของการทดสอบที่ทีมจัดทำไว้ระบุว่าชุดทดสอบฝั่ง backend และชุดทดสอบส่วนประกอบรันไม่ได้ ซึ่งเป็นปัญหาของเครื่องที่ใช้จัดทำเอกสารนั้น ไม่ใช่ปัญหาของโค้ด เพราะทั้งสองชุดผ่านทั้งหมดบนเครื่องที่มีฐานข้อมูลจริง ข้อที่สอง รายการดังกล่าวจัดทำวันที่ 12 กันยายน จึงยังระบุว่าปุ่มส่งคำขอไม่ได้ต่อกับ backend ซึ่งไม่ตรงกับสภาพปัจจุบันแล้ว

กรณีที่ข้ามตัวเองหนึ่งกรณีเป็นจุดที่ต้องแก้ ชุดทดสอบที่ข้ามตัวเองเมื่อหาองค์ประกอบไม่พบ จะรายงานว่า ``ผ่าน'' ในระบบตรวจอัตโนมัติทั้งที่ไม่ได้ตรวจอะไรเลย

% =============================================================================
\clearpage
\section{แผนการดำเนินงานถัดไป}
\label{sec:next}
% ย่อหน้าสุดท้ายยังเกินมาหนึ่งบรรทัด ยืดหน้านี้เล็กน้อยแทนการตัดเนื้อหาเพิ่ม
\enlargethispage{\baselineskip}

\subsection{ครึ่งหลังของ Sprint 4 (18 – 24 กันยายน 2569)}

เนื่องจากจุดวัดข้อ 1 และข้อ 2 ไม่ผ่าน คณะทำงานเรียกใช้แผนสำรองที่ประกาศไว้ในรายงานครั้งที่ 3 คือตัดงานสองรายการออกจาก Sprint 4 เพื่อรักษาเวลาของ Sprint 5 ไว้สำหรับการทดสอบ

\noindent\begin{tabularx}{\dimexpr\textwidth-\leftskip\relax}{@{} L{4.2cm} Y @{}}
\toprule
\hrow \thead{รายการ} & \thead{การตัดสินใจ} \\
\midrule
ตะกร้าข้ามหน่วยงาน & ตัดออกจาก Sprint 4 ตามแผนสำรอง เลื่อนเป็นงานหลังส่งมอบ \\
การจองล่วงหน้า & ตัดออกจาก Sprint 4 ตามแผนสำรอง \\
การอุทธรณ์ & คงไว้ เพราะมีรหัสความต้องการใน SRS และยังไม่มีโค้ดฝั่ง backend \\
การจองห้องและหน้าห้อง & คงไว้ เพราะเป็นขอบเขตของ Sprint 3 ที่ยกมา \\
\bottomrule
\end{tabularx}

งานที่เหลือในครึ่งหลังของ Sprint 4 เรียงตามลำดับที่ต้องทำก่อน คือหน้าส่งมอบของเจ้าหน้าที่ ซึ่งเป็นชิ้นสุดท้ายที่ขวางไม่ให้วงจรยืม–คืนเดินครบผ่านหน้าจอ จากนั้นคือชุดทดสอบผ่านเบราว์เซอร์ที่เดินครบเส้นทาง การเรียกชุดทดสอบนั้นจากระบบตรวจอัตโนมัติ การจองห้องและการอุทธรณ์ฝั่ง backend พร้อมต่อหน้าจอของทั้งสองเรื่อง ปุ่มขอต่ออายุ cron job ที่เหลือ 3 งาน และหน้าส่งออกรายงาน

\subsection{จุดวัดถัดไป}

จุดวัดรอบนี้แยกเป็นข้อละหนึ่งงาน ตามข้อสังเกตในหัวข้อ \ref{sec:review}

\begin{enumerate}
\item หน้าส่งมอบของเจ้าหน้าที่ใช้งานได้จริง ภายใน 19 กันยายน 2569
\item วงจรยืม–คืนเดินครบเส้นทางปกติผ่านหน้าจอ ตั้งแต่ส่งคำขอถึงตรวจสภาพ ภายใน 21 กันยายน 2569
\item ระบบตรวจอัตโนมัติเรียกชุดทดสอบผ่านเบราว์เซอร์ และขวางการรวมโค้ดเมื่อไม่ผ่าน ภายใน 21 กันยายน 2569
\item การจองห้องฝั่ง backend ใช้งานได้ หน้าห้องทั้ง 3 หน้าเลิกใช้ข้อมูลจำลอง และไม่เหลือหน้าว่างในระบบ ภายใน 24 กันยายน 2569 ซึ่งเป็นวันสิ้นสุด Sprint 4
\end{enumerate}

\subsection{การประเมินกำหนดเสร็จ}

คณะทำงานประเมินว่าโครงการยังเสร็จได้ภายใน 1 ตุลาคม 2569 แต่ขอบเขตที่ส่งมอบจะแคบกว่าข้อเสนอโครงการ เพราะตะกร้าข้ามหน่วยงานและการจองล่วงหน้าถูกตัดออกแล้วตามแผนสำรอง

เหตุผลที่ยังคงกำหนดเดิมคืองานที่ค้างมาสองรอบเริ่มเคลื่อนแล้ว และเป็นงานประเภทเดียวกับที่เหลือ คือการต่อหน้าจอเข้ากับ procedure ที่มีอยู่แล้ว ซึ่งประเมินเวลาได้

เหตุผลที่ความมั่นใจลดลงมีสองข้อ ข้อแรก ความสามารถใหม่ฝั่ง backend หยุดนิ่งทั้งรอบ ขณะที่ยังเหลือการอุทธรณ์และการจองห้องซึ่งยังไม่มีโค้ดเลย และทั้งสองต้องเสร็จก่อน Sprint 5 ซึ่งไม่รับงานเพิ่มความสามารถ ข้อที่สอง ระบบตรวจอัตโนมัติยังไม่เรียกชุดทดสอบผ่านเบราว์เซอร์ ความถดถอยของเส้นทางหลักจึงไม่ถูกจับจนกว่าจะมีคนรันเอง

% =============================================================================
\clearpage
\section*{ภาคผนวก — ค่าที่ใช้พล็อตกราฟ}

ค่าเป้าหมายรายมิติ ณ วันสิ้นสุดของแต่ละ Sprint คูณด้วยน้ำหนักของมิตินั้นแล้วรวมกัน คอลัมน์ ``รวม'' คือค่าที่พล็อตบนเส้นแผนในรูปที่ \ref{fig:scurve} ตารางนี้เป็นตารางเดียวกับรายงานครั้งที่ 2 และครั้งที่ 3 ไม่มีการปรับแผน

\noindent\begin{tabularx}{\dimexpr\textwidth-\leftskip\relax}{@{} L{3.2cm} L{1.0cm} L{1.0cm} L{1.0cm} L{1.0cm} L{1.0cm} L{1.0cm} Y @{}}
\toprule
\hrow \thead{ขอบ Sprint} & \thead{BE Des} & \thead{FE Des} & \thead{BE Dev} & \thead{FE Dev} & \thead{Test} & \thead{Doc} & \thead{รวม} \\
\hrow \thead{} & \thead{10\%} & \thead{10\%} & \thead{30\%} & \thead{25\%} & \thead{15\%} & \thead{10\%} & \thead{100\%} \\
\midrule
เริ่ม 24 ก.ค. & 0 & 0 & 0 & 0 & 0 & 10 & 1 \\
S0 · 30 ก.ค. & 45 & 30 & 5 & 5 & 0 & 20 & 12 \\
S1 · 13 ส.ค. & 88 & 75 & 25 & 42 & 10 & 65 & 42 \\
S2 · 27 ส.ค. & 95 & 85 & 60 & 65 & 35 & 70 & 65 \\
S3 · 10 ก.ย. & 100 & 92 & 78 & 78 & 52 & 76 & 78 \\
$*$ กึ่งกลาง S4 · 17 ก.ย. & 100 & 96 & 88 & 88 & 66 & 82 & 86 \\
S4 · 24 ก.ย. & 100 & 100 & 97 & 97 & 80 & 88 & 94 \\
S5 · 1 ต.ค. & 100 & 100 & 100 & 100 & 100 & 100 & 100 \\
\bottomrule
\end{tabularx}

แถวที่มีเครื่องหมายดอกจันคือวันรายงานรอบนี้ ซึ่งไม่ใช่ขอบ Sprint แต่เป็นวันที่ 7 จาก 14 ของ Sprint 4 พอดี ค่าทุกช่องในแถวนั้นจึงเป็นค่ากึ่งกลางระหว่างแถว S3 กับแถว S4 ไม่ใช่ค่าที่กำหนดขึ้นใหม่ วิธีสอดแทรกนี้เป็นวิธีเดียวกับที่รายงานครั้งที่ 2 ใช้

ค่าผลจริง ณ วันรายงาน คือ 97, 96, 79, 80, 47 และ 85 ตามลำดับมิติเดียวกัน ถ่วงน้ำหนักแล้วได้ 78.55 ปัดเป็น \actualpct

ค่าผลจริงรอบก่อนหน้าที่อยู่บนเส้นผลจริง คือ 52\% ณ 27 สิงหาคม 70\% ณ 4 กันยายน และ 75\% ณ 10 กันยายน ตามที่ส่งในรายงานครั้งที่ 1 ถึงครั้งที่ 3 ไม่มีการแก้ย้อนหลัง

\end{document}
